\documentclass{article}
\usepackage[utf8]{inputenc}
\usepackage{amsmath}
\usepackage{geometry}
\geometry{margin=2cm}
\begin{document}

\section*{Análise da Questão 140 — ENEM 2024 — Matemática}

\section{Solução Técnica}

A solução consiste em calcular o número de acertos de João e Felipe com base nas informações fornecidas, e então converter as porcentagens obtidas nos conceitos indicados.

\subsection{Estratégia}
Calcular o número de acertos de:
\begin{itemize}
  \item João: usando 75\% das 200 questões.
  \item Felipe: acertou 30\% a menos que João, ou seja, 30\% a menos da quantidade de questões que João acertou.
\end{itemize}
Determinar a porcentagem de acertos de Felipe e, por fim, identificar os conceitos para ambos conforme a tabela fornecida.

\subsection{Passos}
\begin{enumerate}
  \item Calcular acertos de João:
  \[
    0,75 \times 200 = 150
  \]
  \item Calcular acertos de Felipe:
  \[
    150 - 0,30 \times 150 = 150 - 45 = 105
  \]
  \item Porcentagem de acertos de Felipe:
  \[
    \frac{105}{200} \times 100\% = 52{,}5\%
  \]
  \item Classificar os conceitos:
  \begin{itemize}
    \item João: 75\% está na faixa \(75 \leq P < 90\), conceito \textbf{muito bom};
    \item Felipe: 52{,}5\% está na faixa \(50 \leq P < 60\), conceito \textbf{regular}.
  \end{itemize}
\end{enumerate}

\subsection{Fórmulas Utilizadas}
\begin{align*}
\text{Número de acertos} &= \text{porcentagem de acertos} \times \text{total de questões} \\
\text{Porcentagem} &= \left( \frac{\text{número de acertos}}{\text{total de questões}} \right) \times 100\%
\end{align*}

\subsection{Conclusão}
Os conceitos respectivos de João e Felipe são \textbf{muito bom} e \textbf{regular}.


\section{Explicação Didática}

\subsection{Resumo}
Para determinar os conceitos atribuídos a João e Felipe, calculamos quantas questões cada um acertou a partir das informações porcentuais e depois classificamos seus desempenhos segundo os intervalos indicados.

\subsection{Passo a passo}
\begin{enumerate}
  \item João acertou 75\% das 200 questões:
  \[
    0{,}75 \times 200 = 150
  \]
  \item Felipe acertou 30\% a menos do que João, ou seja:
  \[
    150 - 0{,}30 \times 150 = 150 - 45 = 105
  \]
  \item Determinar a porcentagem de Felipe:
  \[
    \frac{105}{200} \times 100\% = 52{,}5\%
  \]
  \item Classificar segundo os conceitos:
  \begin{itemize}
    \item João: 75\% está no intervalo $75 \leq P < 90$, conceito \'muito bom\'.
    \item Felipe: 52{,}5\% está no intervalo $50 \leq P < 60$, conceito \'regular\'.
  \end{itemize}
  \item Portanto, os conceitos são respectivamente: \textbf{João - muito bom;} \textbf{Felipe - regular}.
\end{enumerate}

\subsection{Erros comuns}
\begin{itemize}
  \item Interpretar "30\% a menos" como 75\% - 30\% = 45\%, em vez de calcular 30\% da quantidade e subtrair.
  \item Confundir os limites dos intervalos de conceito, classificando 75\% em "bom" ao invés de "muito bom".
  \item Calcular incorretamente o percentual de acertos de Felipe após a subtração da quantidade de questões.
  \item Arredondar imprecisamente o valor percentual de Felipe e escolher o conceito errado baseado nisso.
\end{itemize}

\subsection{Dicas para ensino}
\begin{itemize}
  \item Explique a distinção entre reduzir uma quantidade em 30\% e subtrair 30 pontos percentuais de uma porcentagem.
  \item Utilize exemplos práticos para ilustrar o cálculo de porcentagem sobre quantidades e o cálculo do percentual.
  \item Reforce como trabalhar corretamente com intervalos numéricos e classificações baseadas em faixas.
  \item Incentive os alunos a fazerem e conferirem os cálculos passo a passo.
\end{itemize}

\section{Distratores}
\begin{itemize}
  \item \textbf{Alternativa A}: Aplica a redução de 30\% diretamente na porcentagem de João, ou seja, 75\% - 30\% = 45\%, classificando Felipe erroneamente.
  \item \textbf{Alternativa C}: Supõe que o percentual de Felipe é inferior a 50\%, classificando-o como "insatisfatório", devido a erro no cálculo ou arredondamento.
  \item \textbf{Alternativa D}: Classifica João com conceito "bom" (60 \leq P < 75) apesar de ter exatamente 75\%, que pertence à categoria "muito bom", erro comum na interpretação dos limites.
  \item \textbf{Alternativa E}: Coloca João como "bom" e Felipe como "insatisfatório" devido a erros em ambos os cálculos e classificações.
\end{itemize}

Esses distratores são plausíveis porque apresentam erros comuns na interpretação do enunciado e nos intervalos, mas estão incorretos conforme a análise detalhada.

\section{Perfil da Questão}

\begin{itemize}
  \item \textbf{Habilidade Primária}: Interpretação e aplicação de porcentagem.
  \item \textbf{Habilidades Secundárias}: Cálculo de percentual; classificação por intervalos numéricos.
  \item \textbf{Demanda Cognitiva}: Média.
  \item \textbf{Dificuldade Estimada}: Média.
  \item \textbf{Requisitos Específicos}:
  \begin{itemize}
    \item Não há necessidade de interpretação visual adicional.
    \item Requer cálculos numéricos de porcentagem.
    \item Não exige modelagem complexa.
  \end{itemize}
  \item \textbf{Notas}: 
  \begin{itemize}
    \item Fundamental entender "30\% a menos da quantidade" versus "subtrair 30 pontos percentuais".
    \item Conhecimento dos intervalos para conceitos é essencial.
    \item Avalia habilidade de raciocínio matemático contextualizado.
  \end{itemize}
\end{itemize}



\end{document}